\cotefile{79}
\scanpage{65}
\deleted{\unsure{le}} On voit que $G_M^{\unsure{0|\circ}}$ est transitif sur \deleted{toutes} \deleted{\unsure{des|les}} ens. $A_{M,i}$ $(i\in\pi_M)$, et \deleted{\unsure{même}} \deleted{simplement} transitif sur l'ens. $\unsure{\overrightarrow A|\overline A}_{M,i}$ des arêtes orientées ou repères du graphe $\deleted{S_M,A_{M,i}}\,\Sigma(i)$ ;
\deleted{la} transitivité provient du fait que si $(x,y),(x',y')\in\operatorname{Rep}(M)$ sont tels que $x\wedge y=\pm x'\wedge y'$, alors l'unique $u\in\mathrm{Gl}_{\Lambda}(M)$ tel que $(x',y')=u.(x\unsure{.|,}y)$ est dans $\mathrm{Gl}_{\Lambda}(M)^{\pm1}$, donc $\bar u\in G_M^{\unsure{\prime|{}}}$ est dans $G_M^0$.
[La \uline{simple} transitivité est fausse, sauf si $2\,1_{\Lambda}=0$ donc $S_M\simeq M^{**}$, $A_M=$ ens des bases non ordonnées $\{x,y\}$ de $M$\deleted{\unsure{, de}} ; cf plus bas]
\deleted{provient du fait que \unsure{qu'on ait pour} $u.(x,y)\in\{\pm(x,y)\}$}
\deleted{il f et s. que l'on ait $\bar u\in\{\pm\mathrm{id}_M\}$ i.e. $\bar u=1_{G_M}$.}

Il en résulte que le groupe $G_M$ est transitif sur $A_M$, et \deleted{\unsure{même}} \deleted{simplement} transitif sur l'ens $\overrightarrow A_M$.

Introduisons aussi \deleted{\unsure{F}}
\[
\begin{aligned}
 F_M=\bigl\{\deleted{\unsure{(x,y,z)}}f\in\mathfrak P_3(S_M)\mid{}&
 \exists\,x,y,z\in\unsure{B^{*}(M)|B^{*}M)|\widetilde S M}\\
 &\bigl(\deleted{\unsure{x+y}}\;\substack{x+y+z=0\\x\wedge y\in(\det M)^*}\bigr)\\
 &\text{tels que } f\deleted{\unsure{,}}=\{\dot x,\dot y,\dot z\}\bigr\}
\end{aligned}\tag*{(157).}
\]
Alors
\[
\begin{aligned}
\deleted{\unsure{(158)}\ \forall}\quad
\deleted{\unsure{\forall f\in F_M}}\quad
\left\{\begin{aligned}
&\forall f\in F_M,\ \forall a\in\mathfrak P_2(f)\ \deleted{\unsure{\Rightarrow}},\
 \deleted{\unsure{a\in}}\ \text{on a }a\in A_M\\
&\deleted{\unsure{f\in F_M,\ a\in f}}\ \text{de plus, }\tau(a)\text{ ne dépend que de }f,\text{ soit }=\tau(f)\\
&\hfill \unsure{\in|\rotatebox{90}{$\in$}}\ \pi_M .
\end{aligned}\right.
\end{aligned}\tag*{(168)}
\]
% archive sheet number: 64
\scanpage{66}\mspage{32}
donc on a une application
\[
 F_M\xrightarrow{\tau_F}\unsure{\pi|\Pi}\,.\tag*{(169)}
\]
Par ailleurs, on a à distinguer deux cas
\[
 A)\quad 2.1_\Lambda=0\qquad\text{alors }\forall a\in A_M\quad\exists\unsure{!|1} f\ni F_M,\ f\supset a\tag*{(170)}
\]
\deleted{d'où} d'où
\diagram{\begin{tikzcd}[ampersand replacement=\&,column sep=large]
{\text{(\unsure{70|170} bis)}} \& A_M \arrow[r,swap,"{\substack{\text{hom.}\\\text{surjectif}\\\text{de degré }3}}"] \arrow[rr,bend right=45,swap,"\tau_A"] \& F_M \arrow[r,"\tau_F","\text{comm.}"'] \& \pi\,, 
\end{tikzcd}}
et le cas
\[
\begin{aligned} B)\quad 2.1_\Lambda\ne0\qquad &\text{alors }\forall a\in A_M,\text{ il existe}\\
 &\text{exactement \uline{deux} }f\ni F_M\ \text{tels que }f\supset a
\end{aligned}\tag*{(171)}
\]
Dans le cas A), on a
\[
 \overrightarrow A_M\xleftarrow{\sim}\underbrace{\operatorname{Rep}(M)/{\pm\mathrm{id}_M}}_{\operatorname{Rep}'(M)}
 \quad(\text{car }2.1_\Lambda=0)\tag*{(172)}
\]
% slanted
\margin{donc $G_M$ est simplement transitif sur $\overrightarrow A_M$, i.e. $\overrightarrow A_M$ est un $G_M$-torseur.}
Dans le cas B), l'application canonique
\[
 \left\{\begin{aligned}\operatorname{Rep}'(M)&\longrightarrow\overrightarrow A_M\\
 (x,y)^{\bullet}&\longmapsto(\dot x,\dot y)\end{aligned}\right.
 \quad\begin{gathered}\text{est de degré}\\\text{deux,}\\(2.1_\Lambda\ne0)\end{gathered}\tag*{(173)}
\]
elle se factorise en
\[
\begin{aligned}
 \operatorname{Rep}'(M)&\xrightarrow{\sim}\operatorname{Rep}(P_M)\overset{\mathrm{def}}{=}
 \bigl\{(s,t,u)\in S_M\mid\\
 &\qquad\{s,t,u\}\in F_M\bigr\}\\
 \deleted{\unsure{\operatorname{Rep}'(M)|\operatorname{Rep}(M)}}\ (x,y)^{\bullet}&\longmapsto
 \bigl(\dot x,\dot y,\unsure{+}(-x-y)^{\bullet}=(x+y)^{\bullet}\bigr)
\end{aligned}\tag*{(174)}
\]
suivie de
\[
 \begin{aligned}\operatorname{Rep}(P_M)&\longrightarrow\overrightarrow A_M\\
 (s,t,u)&\longmapsto(s,t)\end{aligned}\quad\bigg|\quad
 \begin{gathered}\text{qui est \inserted{bien} de}\\\text{degré }2,\\\text{par (171).}\end{gathered}\tag*{(175)}
\]
% archive sheet number: 65
\scanpage{67}
Dans le cas A, on appellera \uline{repère} \uline{du} graphe \inserted{$P_M=$} $(\unsure{S|B}_M,A_M)$ un élément de $\overrightarrow A_M$, \unsure{ie.|i.e.}
\[
 \operatorname{Rep}(P_M)\overset{\mathrm{def}}{=}\overrightarrow A_M\quad(\text{cas }2.1_{\Lambda}=0),\tag*{(176)}
\]
dans le cas B), on définit $\operatorname{Rep}(P_M)$ \deleted{\unsure{(174)}} par (174). Dans tous les cas, l'ens des repères \unsure{de.|de|abs.} $P_M$ est un torseur sous $G_M$.

\uline{Relations entre les données} $\deleted{S}A_M,F_M$. Je voudrais m'assurer que ces données sur $S_M$
\[
 \left[\begin{gathered}A_M\subset\mathfrak P_2(S_M)\\A_M\longrightarrow\unsure{\pi|\Pi}\end{gathered}\right],\
 \left[\begin{gathered}F_M\subset\mathfrak P_3(S_M)\\F_M\longrightarrow\pi\end{gathered}\right]\tag*{(177)}
\]
% slanted
\margin{$A_M\mathbin{-}\pi$, $F_M\mathbin{-}\pi$ sont \unsure{interprétés|interprétées} ici comme des partitions de $A_M,F_M$ --- \unsure{le|la} \deleted{\unsure{but}} de \unsure{l'ens d'indices|l'un d'eux} $\pi$ étant pour le moment, \unsure{irrelevant}\ldots}
se déterminent mutuellement, En tous cas, dans le cas A\deleted{\unsure{B)} $(2.1_\Lambda\ne0)$}\inserted{comme B}, on a
\[
 A_M=\{a\in\mathfrak P_2(S_M)\mid\exists f\in F_M\deleted{\unsure{\}}},f\supset a\}\tag*{(177)}
\]
(168, \inserted{170,} 171), et d'autre part je \deleted{\unsure{dis}} \inserted{voudrais} \deleted{que}\inserted{\uuline{ds le cas} B $(2.1_\Lambda\ne0)$}
\[
 F_M=\{f\in\mathfrak P_3(S_M)\mid\forall a\in\mathfrak P_2(f),\text{ on a }a\in A_M\}\tag*{\deleted{(178)}}
\]
% slanted
\margin{relation fausse !!}
On a bien $F_M\subset\{\ \}$, par (168), inversement \deleted{\unsure{si on|a-t-on}} a-t-on $F_M\supset\{\ \}$ ?

\uuline{C'est ultra faux} même sur $\Lambda=\mathbb Z/n\mathbb Z$ ou $\Lambda=\mathbb Z$, cela signifierait que si $x,y,z\in M$ sont tels que $x\wedge y,y\wedge z,z\wedge x$ sont dans $\det M^*$, alors on peut \deleted{\unsure{que}} trouver
% archive sheet number: 66
\scanpage{68}\mspage{33}
$\varepsilon,\varepsilon',\varepsilon''\in\{\pm1\}$ tels que l'on ait $\varepsilon x+\varepsilon'y+\varepsilon''z=0$ ou encore qu'on a une égalité de la forme
\[
 z=\deleted{\unsure{\pm}}\varepsilon x+\varepsilon'y\qquad(\varepsilon,\varepsilon'\in\{\pm1_\Lambda\})\tag*{\deleted{(179)} $(*)$}
\]
ce qui est évidemment faux. Par contre, \deleted{\unsure{pour tout $i\in$|pour tout si $i\in$}} \inserted{si on se donne}, en plus de $A_M$, l'application $\tau_A:A_M\longrightarrow\unsure{\pi|\Pi}_M$, i.e la décomposition de $A_M$ suivant les types $i\in\pi$, on réconstitue $F_M$ comme
\[
 \left\{\begin{aligned}
 F_M&=\coprod_{i\in\pi}F_{M,i}\\
 F_{M,i}&=\{f\in\mathfrak P_3(S_M)\mid\forall a\in\mathfrak P_2(f),\text{ on a }a\in A_{M,i}\}
 \end{aligned}\right.\tag*{(\unsure{178|180})}
\]
cette dernière relation en effet signifie que si $x,y,z\in M$ sont tels que
\[
 x\wedge y,\ y\wedge z,\ z\wedge x\quad\in i=\{e,-e\}\subset(\det M)^*
\]
alors on doit avoir une relation $(*)$ ci dessus, en effet. \unsure{comme|Si} $(x,y)\in\operatorname{Rep}(M)$, on aura $z=\varepsilon x+\varepsilon'y$, \inserted{\unsure{avec}} $\varepsilon,\varepsilon'\in\Lambda$, et on aura, en multipliant cette relation par $x$
\[
 \varepsilon'x\wedge y=x\wedge z\qquad\text{d'où}\quad \varepsilon'\in\{\pm1_\Lambda\}
\]
et de \unsure{$=$|même} $\varepsilon\in\{\pm1_\Lambda\}$, cqfd.

Dans le cas \deleted{\unsure{B}}A) par contre $(2.1_\Lambda=0)$,
% archive sheet number: 67
\scanpage{69}
il semble\inserted{rait} que la structure \unsure{$S$}
\[
 \left\{\begin{aligned}F_M&\subset\mathfrak P_3(S_M)\\ F_M&\xrightarrow{\tau_F}\unsure{\pi|\Pi}_M\end{aligned}\right.
\]
soit nettement plus fine que la structure $A_M\subset\mathfrak P_2(S_M),\ A_M\xrightarrow{\tau_A}\pi_M$ qui s'en déduit, via (177) (168) :
\deleted{Une application de $\bar M$ \unsure{supposons}}
$\Lambda$ un corps $k$ de car. 2, \deleted{\unsure{$\ne$}} corps premier $\mathbb F_2$, alors une \deleted{bijection}\inserted{permutation} de \deleted{\unsure{$M$}}\inserted{$M^*=M\smallsetminus\{0\}$ (où $M=k^2$)} peut fort bien transformer couples de pts \deleted{\unsure{qui sont vus en|qui sont une base}} \deleted{\unsure{par}} \deleted{alignés}\inserted{colinéaires} en couples \deleted{alignés}\inserted{colinéaires} --- donc aussi bases en bases --- et, je suppose, de façon compatible avec les jauges i.e. en transformant \deleted{\unsure{base}} deux bases de \unsure{$=$|même} jauge en \deleted{deux} \inserted{des} bases de $=$ jauge, sans pour autant que \unsure{respecte|respecter} la relation $x+y+z=0$, pour \deleted{\unsure{$x,y,z$}}, (\inserted{avec} $x\wedge y=y\wedge z=z\wedge x\in(\det M)^*$) \ldots.
% slanted
\margin{\uline{à élucider} éventuellement, en utilisant la structure d'addition non partout définie de $\pi_M=(\det M)^*$, \deleted{\unsure{en}} et exigeant qu'elle soit respectée par les \deleted{\unsure{jauges}} \unsure{autom|permutations}, on arrive à déduire parfois $F_M$ de $A_M,\ \tau_A:A_M\to\pi_M$ \ldots}

À vrai dire, la structure que nous étudions en car. 2, se réduisant \inserted{alors} à l'ens. des sommets $\widetilde S_M=M^{**}$ (ens. des éléments de $M$ faisant partie d'une base), l'ens. $\widetilde F_M\unsure{\subset|\not\subset|\notin}\mathfrak P_3(S)$, formé de parties \inserted{$\{x,y,z\}$} de $\widetilde S_M$ satisfaisant $x+y+z=0$, $x\wedge y\ (=y\wedge z=z\wedge x)\in(\det M)^*$, et de l'application
\[
 \widetilde\tau:\widetilde F_M\longrightarrow(\det M)^*=\widetilde\pi\quad
 \{x,y,z\}\longmapsto\begin{aligned}[t]x\wedge y&=y\wedge z\\&=z\wedge x\end{aligned}\tag*{(178)}
\]
% archive sheet number: 68
\scanpage{70}\mspage{34}
peut avoir un intérêt sur un anneau quelconque $\Lambda$, pas néc. de car. 2 i.e. tel que $2\,1_\Lambda=0$, à cela près que l'application $\widetilde\tau$ est alors à valeurs non dans $\widetilde{\unsure{\pi|\Pi}}$, mais dans $\pi=\widetilde\pi/\{\pm1_\Lambda\}$
\[
\begin{aligned}
 \widetilde\tau:\widetilde F_M&\longrightarrow\pi=(\det M^{\unsure{*|\wedge}})/\{\pm1\}\\
 \{x,y,z\}&\longmapsto x\wedge y\bmod\pm1
\end{aligned}\tag*{(178 bis)}
\]
--- si on veut une application à valeurs dans $\det M^*$, il faudra prendre
\[
\begin{aligned}
 \widetilde\tau_{\mathrm{or}}:\underbrace{\widetilde F_{.\mathrm{or}.M}}_{\text{``faces orientées''}}&
 \longrightarrow\widetilde\pi=\det M^*\\
 \left.\begin{gathered}(x,y,z)\\(y,z,x)\\(z,x,y)\end{gathered}\right|&\longmapsto x\wedge y=y\wedge z=z\wedge x
\end{aligned}\tag*{(179)}
\]
Notons qu'on a même une application analogue sur $F_{.\mathrm{or}.M}\simeq\widetilde F_{.\mathrm{or}\,M}/{\pm\mathrm{id}_M}$, puisque l'on a
\[
 \widetilde\tau_{\mathrm{or}}\deleted{\unsure{(x,y,z)}}\inserted{(\overrightarrow{\widetilde f})}
 =\widetilde\tau_{\mathrm{or}}(-\mathrm{id}_M(\overrightarrow{\widetilde f}))
\]
pour une face orientée de $P\widetilde M$ --- d'où
\[
 F_{.\mathrm{or}\,M}\xrightarrow{\tau_{\mathrm{or}}}\deleted{\unsure{\det}}\widetilde\pi\tag*{(180)}
\]
qui \deleted{\unsure{relie}} s'insère dans le \unsure{diagr.|diag.} comm.
\diagram{\begin{tikzcd}[ampersand replacement=\&,column sep=large]
{\text{(181)}} \& F_{.\mathrm{or}.M}\arrow[r]\arrow[d,"\text{can}"'] \& \widetilde\pi\arrow[d]\\
\& F_M\arrow[r] \& \pi\,,
\end{tikzcd}}
% archive sheet number: 69
\scanpage{71}
et la question se pose dans quelle mesure cette structure supplémentaire se déduit des autres structures \ldots.

\deleted{\uline{\unsure{Récapitulation}}} \deleted{\unsure{Soit}}
% archive sheet number: 70
\scanpage{72}\mspage{35}
$\overline{\underline{\mathrm V}}$\quad Récapitulation

Soit $\Lambda$ un anneau comm. non nul, $M$ un $\Lambda$-module libre de rang 2. On lui associe deux ``structures \unsure{triangulées|triangulières}'' (ens. de sommets $S$, plus \unsure{une|un} \inserted{ens. de} parties \deleted{\unsure{à}} \inserted{$F$} $\unsure{\subset|\in}\mathfrak P_3(S)$ \deleted{\unsure{ayant comme}} \inserted{de} réunion dans $S$), dont la deuxième n'est intéressante que si $2.1_\Lambda\ne0$ (autrement, elle coïncide avec la première).

Les structures triangulées \deleted{sont} en question sont ``partitionnées'' (i.e. $F\subset\mathfrak P_3(S)$ est munie d'une partition, via une application sur un ens d'indices convenable $I$) --- donc la structure envisagée peut être considérée comme un ens. de structures triangulées sur le \unsure{$=$|même} ens $S$, indexé par $I$. Les deux structures envisagées dépendent fonctoriellement de $M$, et la deuxième est le quotient de la première par $\{\pm\mathrm{id}_M\}\subset\mathrm{Gl}_\Lambda(M)$.

A)
\[
 \left\{\begin{aligned}
 \widetilde S_M&=M^{*\unsure{*}}=\{x\in M\mid\exists y\in M,\ (x,y)\in\operatorname{Rep}(M)\}\\
 \widetilde F_M&=\deleted{\unsure{\mathfrak P}}\bigl\{\{x,y,z\}\in\mathfrak P_3(\widetilde S_M)\mid
 \substack{x+y+z=0\\x\wedge y\ (=y\wedge z=z\wedge x)\in(\det M)^*}\bigr\}
 \end{aligned}\right.\tag*{(1)}
\]
\[
 \overrightarrow{\unsure{\pi|\Pi}}_M=\det M^*,\quad \pi_M=(\det M^*)/\{\pm1_\Lambda\}\tag*{(2)}
\]
% archive sheet number: 71
\scanpage{73}
On pose
\[
 \underbrace{\widetilde{\overrightarrow F}_{\deleted{\unsure{\mathrm{or}}}M}}_{\substack{\text{ens. des \uline{faces}}\\\text{\uline{orientées} de }\widetilde P_M=(\widetilde S_M,\widetilde F_M)}}
 =\{(f,\omega)\mid f\in\widetilde F_M,\quad
 \omega\in\underbrace{\underline\omega_f}_{\substack{\text{ens. des}\\\text{ordres circ.}\\\text{sur }f}}\}\tag*{(3)}
\]
\[
 \overrightarrow f\overset{\deleted{\unsure{f}}}{\longmapsto}\check{\overrightarrow f}\quad
 \text{donné par}\quad(f,\omega)\longmapsto(f,\omega^{-1})\tag*{(4)}
\]
(NB: on \deleted{\unsure{f}} réserve la notation $-\overrightarrow f$ à la face orientée déduite de $\overrightarrow f$ par transport de structure par $-\mathrm{id}_M$)
\[
 \begin{aligned}
 \widetilde{\overrightarrow F}_{\deleted{\unsure{\mathrm{or}}}M}&\xrightarrow{\widetilde{\overrightarrow\tau}_{\mathrm{or}}}\overrightarrow{\unsure{\pi|\Pi}}_M\\
 \widetilde\tau(\overrightarrow f_{x,y,z})&=x\wedge y\quad(=y\wedge z=z\wedge x)
 \end{aligned}\tag*{(5)}
\]
% slanted
\margin{on note $\overrightarrow f_{x,y,z}\in\widetilde F_{\mathrm{or}\,M}$ la face orientée $(\{x,y,z\},\unsure{x|a}\to y\to z\to x)$}
On a
\[
 \widetilde{\overrightarrow\tau}_{\mathrm{or}}(\check{\overrightarrow f})=
 -\widetilde{\overrightarrow\tau}_{\mathrm{or}}(\overrightarrow f)
\]
d'où par passage aux quotients une appl.
\[
 \widetilde F_M\xrightarrow{\widetilde\tau}\pi_M\,.\tag*{(6)}
\]
Pour $i\in\pi$, on pose
\[
 \left\{\begin{aligned}
 \widetilde F_{M,i}&=\widetilde\tau^{-1}(\{\unsure{{}'i|i}\})\qquad\text{d'où}\quad\widetilde F_M=\coprod_{i\in\pi}\widetilde F_{M,i}\\
 \widetilde P_{\deleted{\unsure{M}}M,i}&\text{ est le ``polyèdre triangulé'' }(\widetilde S_M,\widetilde F_{M,i})
 \end{aligned}\right.\tag*{(7)}
\]
\deleted{Proposition} On pose
\[
 \begin{aligned}
 \underset{\text{\uline{arêtes de} }\widetilde P_M}{\widetilde A_M}
 &=\{a\in\mathfrak P_2(\widetilde S_M)\mid\exists f\in\widetilde{\deleted{S} F}_M,\ f\supset a\}\\
 &\bigl[=\{\{x,y\}\in\mathfrak P_2(\widetilde S_M)\mid(x,y)\in\operatorname{Rep}(M)\}\bigr].
 \end{aligned}\tag*{(8)}
\]
et on a
% archive sheet number: 72
\scanpage{74}\mspage{36}
% archive sheet number: 73
\deleted{\uline{Proposition}}\inserted{(9)} Toute arête de $\widetilde P_M$ est contenue dans \deleted{une unique face} \inserted{une unique face}, d'où
\[
 \widetilde A_M\xrightarrow[\text{degré }3]{}\widetilde F_M
\]
et une application composée
\diagram{\begin{tikzcd}[ampersand replacement=\&,column sep=large]
{\text{(10)}} \&\widetilde A_M\arrow[r]\arrow[rr,bend right=22,"\widetilde\tau_A"'] \&\widetilde F_M\arrow[r,"\widetilde\tau_F"] \&\unsure{\pi|\Pi}_M
\end{tikzcd}}
De \unsure{$=$|même}, avec la définition évidente des arêtes orientées, on a
\diagram{\begin{tikzcd}[ampersand replacement=\&,column sep=large]
{\text{(11)}} \&{\underset{\substack{\text{arêtes orientées}\\\text{de }\widetilde A_M}}{\widetilde{\overrightarrow A}_M}}\arrow[r,"\text{degré }3"]\arrow[rr,bend right=31,"\widetilde{\overrightarrow\tau}_A"'] \&\widetilde{\overrightarrow F}_M\arrow[r,"\widetilde{\overrightarrow\tau}_F"] \&\overrightarrow\pi_M
\end{tikzcd}}
\margin{\uline{NB $\widetilde{\overrightarrow\tau}_A(\check{\overrightarrow a})=-\widetilde{\overrightarrow\tau}_A(\overrightarrow a)$}}
% slanted
\margin{(10) \deleted{\unsure{peut}} être vu comme déduit de (11) en passant aux quotients par $\overrightarrow a\longmapsto\check{\overrightarrow a}$}
On définit $\widetilde A_{M,i}$ \deleted{\unsure{en}} (pour $i\in\pi_M$) comme l'image inverse de $\{i\}$ par $\widetilde{\overrightarrow\tau}_A$, d'où
\[
 \widetilde A_M=\coprod_{i\in\pi}\widetilde A_{M,i}\,.\tag*{(12)}
\]
\[
\boxed{\begin{gathered}
 \deleted{\widetilde A_{M,i}\text{ et }\widetilde F_{M,i}\text{ se déterminent mutuellement}}\\
 \left\{\begin{aligned}
 &\deleted{\widetilde A_{M,i}=\{a\in\mathfrak P_2(\widetilde S_M)\mid\exists f\in\widetilde F_{M,i},f\supset a\}}\\
 &\deleted{\widetilde F_{M,i}=\{f\in\mathfrak P_3(\widetilde S_M)\mid\forall a\in\mathfrak P_2(f),a\in\widetilde A_{M,i}\}}
 \end{aligned}\right.
\end{gathered}}\tag*{\deleted{(13)}}
\]
\deleted{\unsure{Dans}} On peut considérer que $(\widetilde F_{M,i})_{i\in\pi_M}$ définit une famille de structures triangulées sur le $=$ ens. de sommets $\widetilde S_M$, paramétré par $\pi_M$.
% slanted
\margin{\deleted{\unsure{ou encore, on a} une structure}}
Il faudra à l'occasion interpréter la structure
% slanted
\margin{\uline{tirer au} \uline{clair}}
\[
 \pi_M\text{ torseur sous }\deleted{\unsure{\Lambda}} Z'_{\Lambda}=\Lambda^*/\{\pm1_\Lambda\},\tag*{(1\deleted{\unsure{5}}3)}
\]
et aussi
\[
 \overrightarrow\pi_M\text{ torseur sous }\Lambda^*\text{ qui ``relève'' }\pi_M\,.\tag*{(1\deleted{\unsure{6}}4)}
\]
\scanpage{75}
\uline{Remarque} Il \uline{n'est pas vrai} que la seule donnée des $\widetilde A_{M,i}$ permet de retrouver \deleted{les} $\widetilde F_{M,i}$, par la \unsure{récette|recette} naïve --- car \inserted{pour} un $f\in\mathfrak P_3(\widetilde S_M)$, $f=\{x,y,z\}$, il ne suffit pas qu'on ait $a\in\widetilde A_{M,i}\ \forall a\in\mathfrak P_2(f)$, i.e $x\wedge y,y\wedge z,z\wedge x\in i=\{\pm e\}\subset\det M^*$, pour avoir $f\in\widetilde F_{M,i}$ i.e. $x+y+z=0$. En effet, on constate que la condition envisagée signifie simplement que $z=\pm x\pm y$ (et non $z=-(x+y)$.
Mais dans le cas $2.1_\Lambda=0$ --- qui est celui où c'est la structure $\widetilde P_M$ \deleted{\unsure{qui}}\inserted{elle-même} \unsure{est|où} la plus intéressante --- cela implique bien $x+y+z=0$, i.e.
\[
 \text{Si }2.1_\Lambda=0,\quad\text{alors }\widetilde F_{M,i}=\bigl\{f\in\mathfrak P_3(\widetilde S_M)\mid
 \substack{\forall a\in\mathfrak P_2(f),\\\text{on a }a\in\widetilde A_{M,i}}\bigr\}.\tag*{(15)}
\]

\uline{Question de connexité des} $\widetilde P_{M,i}$. On s'y intéresse, car on aimerait qu'un autom. de $\widetilde P_{M,i}$ (i.e. permutation de $\widetilde S_M$ qui \deleted{\unsure{applique}}\inserted{transforme} $\widetilde F_{M,i}$ en lui-\unsure{$=$|même}) qui fixe un \uline{repère} $(x,y,z)$ (i.e. un $(x,y,z)\in\widetilde S_M^3$ tels que $\{x,y,z\}\in\widetilde F_{M,i}$) soit l'identité.
% slanted
\margin{donc il faut considérer l'antipodisme comme faisant partie de la structure \unsure{enrichie} de $\widetilde P_{M,i}$ \ldots\ '' )!}
[Pour pouvoir travailler, il \unsure{le faut|faut} au moins que l'autom. commute à ``l'antipodisme'' \deleted{\unsure{$\varepsilon$|$\in$}} $\underline{\widetilde a}_M$ (autom de $\widetilde P_{M,i}$ induit par $-\mathrm{id}_M$). Alors, si $u$ fixe un $\unsure{s|\xi|x}\in\widetilde S_M$,
% archive sheet number: 74
\scanpage{76}\mspage{37}
% slanted
\margin{\deleted{\unsure{inutile} ?}}
il fixe aussi $-\deleted{\unsure{x|\xi}}\unsure{s|\xi|x}$.] S'il fixe $\deleted{\unsure{s|\xi}}x$ et $y$ tels que $x\wedge y\in i=\{e,-e\}\subset\deleted{\unsure{\det}}\overrightarrow{\unsure{\pi|\Pi}}\deleted{\unsure{{}^*}}=(\det M)^*$, alors il fixe aussi $\inserted{z=}-(x+y)$, \inserted{donc aussi $x+y$.} Cela nous permet de nous propager à partir d'un \inserted{$i$-}repère $(x,y)$, \deleted{fixé par $u$} \deleted{\unsure{on le transforme soit par $(x,y+x$|les remplaçant par $(x,y+x)$}} i.e. à partir d'une $i$-face $(x,y,z)$, en faisant soit \deleted{des transpositions}\inserted{une permutation} sur $(x,y,z)$ --- p. ex. les deux transpositions entre $x$ et $y$, et entre $y$ et $z$ --- et la transformation $(x,y,z)\longmapsto(x,-y,y-x)$ ---
\margin{\unsure{1}}
en termes matriciels, avec tout \inserted{$i$-}repère \uline{de $M$} $r=(x,y)$, interprété comme isom
\[
 r:\Lambda^2\xrightarrow{\sim}M
\]
qui soit formé de ``composantes'' $x,y$ fixées par $u$ (on dit que $r$ est ``$u$-fixe''), il en est de \unsure{$=$|même} de $r\circ\tau_0^{\unsure{r|v}},\ r\circ\tau_1^r,\ r\circ\tau_\infty^r$ (cf (82) \inserted{p. 15'}) .
Cela signifie que si $r$ est $u$-fixe, de $=$ $r\circ g$, où
\[
 \bigl(\qquad g\in\underset{\substack{\Vert\\\operatorname{Aut}_\Lambda(\Lambda^2)}}{\mathrm{\deleted{\unsure{S}}Gl}(2,\Lambda)}
 \quad\begin{gathered}\text{est dans le sous-groupe}\\\text{engendré par }\tau_0^r,\tau_1^r,\tau_\infty^r.\end{gathered}
\]
Or ce sous-groupe est l'image, par
\[
 \mathrm{Sl}(2,\mathbb Z)\longrightarrow\mathrm{Sl}(2,\Lambda)\subset\mathrm{Gl}(2,\Lambda)
\]
du sous-groupe de $\mathrm{Sl}(2,\mathbb Z)$ engendré par ces $=$s matrices $\tau_i^r$ (interprétées comme étant à
% archive sheet number: 75
\scanpage{77}
coeff. entiers). Comme on sait que ce sous-groupe est $\mathrm{Sl}(2,\mathbb Z)$ tout entier, et que (si $\Lambda_0=\operatorname{Im}(\mathbb Z\mathbin{\unsure{-|\to}}\Lambda)$) \deleted{\unsure{celui}} l'hom
\[
 \mathrm{Sl}(2,\mathbb Z)\longrightarrow
 \underset{\substack{\rotatebox{-90}{$\subset$}\\\mathrm{Sl}(2,\Lambda)\subset\mathrm{Gl}(2,\Lambda)}}{\mathrm{Sl}(2,\Lambda_0)}
 \quad\text{est surjectif}
\]
on voit que le sous-groupe de $\mathrm{\deleted{S}Gl}(2,\Lambda)$ engendré par les $\tau_i^{r}$ n'est autre que \deleted{$\mathrm{Sl}(2,\Lambda_0)$.} \unsure{/}
\[
 \underset{\substack{\rotatebox{-90}{$\subset$}\\\mathrm{Sl}(2,\Lambda)}}{\mathrm{Sl}(2,\Lambda_0)}
 =\left\{\begin{pmatrix}a&b\\c&d\end{pmatrix}\ \middle|\
 \begin{aligned}&a,b,c,d\in\Lambda_0=\operatorname{Im}(\mathbb Z\mathbin{-}\Lambda)\\&ad-bc=1\end{aligned}\right\},
\]
en faisant opérer celui-ci \deleted{à droite} sur l'ens des repères, via
\[
 \deleted{\unsure{(f}}\ (x,y)\longmapsto(ax+by,cx+dy)
\]
Ainsi, on a prouvé

\uline{Lemme} Soit $u$ un autom. de $P_{M,i}$, qui fixe \inserted{deux éléments $x,y\in S_M$ formant} une $i$-base $(x,y)$ de $M$ (i.e. $x\wedge y\in i\subset\overrightarrow{\unsure{\pi|\Pi}}_M=\det M^*$ ---), \deleted{Alors} supposons de plus que $u$ commute à l'antipodisme $-\mathrm{id}_M$. Alors $u$ fixe \deleted{les} tous les éléments de $\widetilde S_M$ qui sont de la forme $ax+by$, avec $a,b\in\Lambda_0=\operatorname{Im}(\mathbb Z\to\Lambda)$, et premiers \deleted{\unsure{ce}} entre eux (i.e. $a\Lambda+b\Lambda=\Lambda$ i.e. $\exists c,d\in\Lambda$ avec $ad-\unsure{cc|ce}=1$).
% archive sheet number: 76
\scanpage{78}\mspage{38}
\uline{Corollaire} Supposons $\Lambda=\Lambda_0$ i.e. $\Lambda$ est un anneau ``absolu''. Alors tout autom de $P_{M,i}$ qui fixe un repère (i.e. les sommets d'une face) et commute à l'antipodisme est l'identité.

On trouve donc

\uline{Théorème} \unsure{Soient|Soit} $\Lambda\simeq\mathbb Z$ ou $\mathbb Z/n\mathbb Z$, \inserted{$M$ un $\Lambda$-mod. libre de rg 2.} Alors

a) Pour tte structure $\tfrac12$-symplectique $i\in\underbrace{(\det M^*)/\{\pm1\}}_{\unsure{\pi|\Pi}_M}$, \deleted{le groupe des automorphismes} \inserted{l'hom. canonique}
\[
 \mathrm{\deleted{\unsure{S}}Gl}_\Lambda(M)^{\pm}\longrightarrow
 \deleted{\unsure{G_{\widetilde S_M}}}\underbrace{\operatorname{Aut}(\widetilde P_{M,i},\deleted{\unsure{\xi}}\mathrm{antipod})}_{\substack{\text{automs qui}\\\text{commutent à}\\\text{l'antipodisme}}}\tag*{(16)}
\]
\deleted{est injectif (sur unité) et} est un isomorphisme [car $\mathrm{Gl}_\Lambda(M)^{\pm}$ est simplement transitif sur l'ens des bases $\tfrac12$ symplectiques relatives à $i$].

b) \deleted{\unsure{Tout autom de}}\inserted{Considérons sur} $\widetilde S_M=\deleted{\unsure{A}}M^*$ la structure $\widetilde P_M$ définie par la famille $(\widetilde F_{M,i})_{i\in\pi_M}$ \deleted{\unsure{est}} de structures triangulées, et l'antipodisme, qui est un autom \inserted{involutif} commun à toutes ces structures triangulées. Alors
% slanted
\margin{il suffit \unsure{$=$|même} de ne regarder \uline{que} la structure donnée par $\widetilde F_M$ tout seul, \unsure{sans} sa partition}
\[
 \mathrm{Gl}_\Lambda(M)\xrightarrow{\sim}\operatorname{Aut}(\widetilde P_M)\,.\tag*{(17)}
\]
% archive sheet number: 77
\scanpage{79}
% archive sheet number: 78
\uline{Scholie} \inserted{Soit $n\in\mathbb N\setminus\{1\}$, d'où $\Lambda_n\simeq\begin{cases}\mathbb Z&\text{si }n=0\\\mathbb Z/n\mathbb Z&\text{sinon}\end{cases}$,} Par les principes généraux, il s'ensuit qu'on a \deleted{\unsure{entre n'importe|réciproquement}} \inserted{\deleted{deux}} \unsure{équivalences|équivalence} de groupoïdes
\[
 (M,i)\longmapsto\widetilde P_{M,i}
\]
\diagram{\begin{tikzcd}[ampersand replacement=\&,column sep=small]
{\text{(18)}} \& {\substack{\unsure{\text{Groupoïde}}\,\inserted{\text{des}}\,\text{modules libres}\\\text{$\frac12$ symplectiques }\deleted{\text{de rang 2}}\inserted{\text{de rang 2}}\text{ sur }\Lambda_n\\\deleted{\unsure{\Lambda_n}}}}\arrow[r,"{\approx}"] \&
{\substack{\text{groupoïde des}\\\text{structures triangulées}\\\text{avec antipodisme}\\\text{(autom involutif)}\\\deleted{\unsure{\text{satisf}}}\text{ isomorphes à}\\\text{une structure du type}\\\text{canonique (dépendant}\\\deleted{\unsure{\text{de }\Lambda}}\text{ de }n\text{)}}}
\end{tikzcd}}
\[
 M\longmapsto\widetilde P_M
\]
\diagram{\begin{tikzcd}[ampersand replacement=\&,column sep=small]
{\text{(19)}} \& {\substack{\text{groupoïde des}\\\text{modules libres de}\\\text{rang 2 sur }\Lambda_n\\\deleted{\unsure{(n=0,\;n\in\mathbb N\setminus\{1\})}}}}\arrow[r,"\approx"] \&
{\substack{\text{groupoïde des}\\\text{structures }(S,F,\unsure{\pi|\Pi},\tau,a)\\\text{où }\deleted{\unsure{(S,F)}}\deleted{\unsure{\text{avec}}}\text{ on a une}\\\deleted{\unsure{F}}F\subset\mathfrak P_3(S)\text{ une str}\\\text{triangulée sur }S,\\\deleted{\unsure{F}}\text{ décomposée en familles}\\\inserted{\text{disjointes}}\text{ indexées par }\pi\\\text{(i.e. on a }\tau:F\longrightarrow\pi\ \text{surj)}\\a\text{ est un autom inv.}\\\text{(antipodisme) de la}\\\text{structure précédente.}}}
\end{tikzcd}}
% slanted
\margin{\uline{NB} Si $2.1_\Lambda=0$ (i.e., ici, $\Lambda=\mathbb F_2$) il est inutile d'introduire l'antipodisme, qui est l'identité. Si $2.1_\Lambda\ne0$ \deleted{\unsure{($\Lambda=\mathbb F_2$ exclu)}} alors l'antipodisme est une involution ss pts fixes...}

Il conviendrait de caractériser axiomatiquement, par des propriétés intrinsèques, les structures qu'on obtient au deuxième mb. De plus, il faudrait comprendre, en termes de ces données, \uline{comment reconstruire} $M$, et par là aussi,
% slanted
\margin{cf suggestion (23) ci-dessous}
\scanpage{80}\mspage{39}
% archive sheet number: 79
les structures supplémentaires reliant les données $S,F$ etc \unsure{à|:} $\Lambda$ (cf. (13)(14)(5) et aussi (118) p. 23)
% slanted
\margin{il y a aussi l'opération de $\Lambda^*$ \deleted{\unsure{par}} sur $\widetilde S_M$}

Enfin, dans le cas où on ne suppose plus $\Lambda=\Lambda_0$, on voudrait introduire sur $\widetilde P_{M,i}$ \inserted{resp. sur} $\widetilde P_M$ un élément de structure supplémentaire (peut-être l'opération de $\Lambda^*$ sur $\widetilde S_M$) qui permette néanmoins d'interpréter $\mathrm{Gl}_\Lambda(M)^{\pm}\deleted{\unsure{(\ )}}$ resp. $\mathrm{Gl}_\Lambda(M)$ comme le groupe d'automorphismes de cette structure.

\deleted{\unsure{Si $2.1_\Lambda\ne0$}} On suppose dorénavant $2.1_\Lambda\ne0$.
\deleted{B) Le groupe $\{\pm\operatorname{id}_M\}$ opère librement sur $\widetilde S_M,\widetilde A_M,\widetilde F_M$ [\uline{NB} si $\{x,y,z\}\deleted{\unsure{\}}}\in\widetilde F_M$, alors les éléments $x,y,z,-x,-y,-z$ sont tous distincts...]. On pose}

Revenons à $\widetilde P_M$ \inserted{dans le cas $\Lambda$ quelconque}, et considérons l'opération de $\Lambda^*$ sur $\widetilde{\deleted{\unsure{B}}S}_M$, par les $\widetilde T_\lambda$
\[
 \widetilde T_\lambda x=\lambda.x\tag*{(20)}
\]
alors $\Lambda^*$ respecte toutes les structures ($\deleted{\unsure{a}}\widetilde F_M\subset\mathfrak P_3(\widetilde S_M)$ et la partition de $\widetilde F_M$ en parties $\widetilde F_{M,i}$) -- l'opération de $\mathrm{Gl}_\Lambda(M)$ sur $\widetilde S_M$ \deleted{est} commute à $\Lambda^*\inserted{\simeq\operatorname{Centre}(\mathrm{Gl}_\Lambda(M))}$, en respectant aussi
% slanted
\margin{\uline{NB} la donnée de l'antipodisme est contenue dans cette opération de $\Lambda^*$}
